\chapter{Conclusion}
\label{ch:conclusion}


\section{Summary of Findings}

This thesis asked one question: in real-time sepsis prediction on MIMIC-IV, how
much of the reported performance comes from the model versus the researcher's
choices about the label, the data split, and the evaluation metric?

The answer is clear: the researcher's choices matter as much as the model, and
the sepsis label itself has a fundamental problem: it contains no signal that
is independent of clinical action.

\paragraph{H1 (Label dominance): Not confirmed under correction.}
Changing the sepsis definition (while keeping the model fixed) changes AUROC by
0.153 (95\% CI 0.136--0.177). Changing the model (while keeping the label
fixed) changes AUROC by 0.142 (95\% CI 0.130--0.161). Both are large relative
to typical performance differences reported in the literature, but the margin
between them is 0.011 (95\% CI $-0.017$--$0.034$, $p = 0.21$, Holm-adjusted
$p = 0.63$), so the label cannot be shown to dominate the model. Cohen et al.'s
\cite{cohen2024subtle} MIMIC-III finding is supported on MIMIC-IV in its weaker
form: label choice contributes variance comparable to model choice.

\paragraph{H2 (Coefficient instability): Confirmed, descriptively.}
Ten of the 20 shared covariates, including MAP, temperature, bilirubin,
platelet count, and vasopressor indicators, reverse sign across label variants;
18 reverse sign or change magnitude at least twofold. H2 is a stability
criterion rather than a contrast, so no test statistic is defined for it.
Published hazard ratios from MIMIC-IV sepsis models are not interpretable
without knowing which label definition was used.

\paragraph{H3 (Treatment leakage): Structural finding.}
The pre-treatment window contains zero sepsis onset events across 179,003
person-hours (exact one-sided 95\% upper bound on the per-hour onset rate:
$1.7\times10^{-5}$). The Sepsis-3 label on MIMIC-IV does not contain any event that
occurs before the clinical action that defines it. Every published MIMIC-IV
real-time sepsis model measures the ability to predict clinician suspicion, not
physiological deterioration.

\paragraph{H4 (Split optimism): Confirmed.}
The random split inflates AUROC by 0.123 (95\% CI 0.100--0.147) compared to the
temporal split for the primary model (0.730 vs.\ 0.607), far exceeding the 0.02
threshold (Holm-adjusted $p = 0.003$). Studies
using random splits substantially overestimate real-world performance.

\paragraph{H5 (Metric divergence): Confirmed.}
The Utility Score is negative for every model and every label variant in the
temporal test set. The utility collapse documented by Wang et al.\
\cite{wang2025srpm} across ninety-one models appears \textit{internally}
on MIMIC-IV under an honest time-based split. External AUROC falls by 0.170
(95\% CI 0.102--0.242, Holm-adjusted $p = 0.004$) to 0.436, an interval whose
upper limit reaches chance. Alert burdens range from 91 to over 340 false
alerts per true alert (benchmark: 1.4, Moor et al.\ \cite{moor2023review}).

\paragraph{H6 (Model-class null): Rejected.}
GBT outperforms the primary model by 0.135 AUROC (0.742 vs.\ 0.607), 95\% CI
0.113--0.156, entirely above the 0.02 equivalence bound (Holm-adjusted
$p = 0.003$). But the clinical ordering reverses: the primary model has a less
negative Utility Score ($-1.004$ vs.\ $-5.059$). The ``better'' model by AUROC
produces more harmful alerts.


\section{The Central Lesson}

The utility collapse documented by Wang et al.\ \cite{wang2025srpm} is
not a model problem. It is a label problem. The Sepsis-3 label on MIMIC-IV:

\begin{enumerate}
    \item \textbf{Depends on clinical action:} no patient is \emph{labelled}
    septic before doctors order antibiotics and cultures. A deterioration label
    that uses no treatment timestamp records 1,362 onsets in the same window, so
    the emptiness belongs to the definition, not to the patients.
    \item \textbf{Is sensitive to how it is defined:} the number of sepsis
    cases nearly doubles across plausible Sepsis-3 interpretations.
    \item \textbf{Produces harmful alert behaviour:} every model, under every
    label variant, has a negative Utility Score.
\end{enumerate}

These three problems are independent of model complexity. A better model
cannot fix a label that is constituted by the very clinical actions it is
supposed to predict ahead of. The solution requires changing the label.

Splitting the definition sharpens point (2) and qualifies it
(Section~\ref{sec:label_anchor_results}). It is qualified because the onset
\emph{time} under every Sepsis-3 variant is fixed entirely by the treatment
anchor --- the organ-dysfunction criterion decides which stays are labelled and
moves no onset by an hour --- so part of the variation between variants is
guaranteed by the definition rather than found in the data. It is sharpened
because a label built from physiology alone, owing nothing to treatment timing,
agrees with Sepsis-3 barely better than chance ($\kappa = 0.057$). The label
effect is therefore real, not merely definitional.

The same comparison relocates the difficulty. Sepsis-3 is \emph{harder} to
predict than either of the limbs it is built from: the treatment anchor alone
reaches 0.756 AUROC and physiological deterioration alone 0.842, against 0.607
for the conjunction. Neither clinician behaviour nor patient physiology is
unreadable. What resists prediction is the requirement that both occur
together, which is what Sepsis-3 operationalises.


\section{Contributions}

\begin{enumerate}
    \item \textbf{C1 (Label variance):} To the best of our knowledge, the first
    three-variant label sensitivity analysis for a dynamic survival model on
    MIMIC-IV, confirming Cohen et al.\ at a larger scale (65,078 stays vs.\
    3,186).

    \item \textbf{C2 (Treatment leakage exposed):} Demonstration that the
    pre-treatment window contains zero onset events, providing the strongest
    available confirmation of Kamran et al.'s \cite{kamran2024evaluation}
    finding.

    \item \textbf{C3 (Utility collapse):} Negative Utility Score for every model
    and label variant, extending Wang et al.'s \cite{wang2025srpm}
    meta-analytic finding to a controlled single-study comparison.

    \item \textbf{C4 (Coefficient instability):} Ten of the 20 shared covariates
    reverse sign across labels and 18 reverse sign or change magnitude at least
    twofold, confirming Lauritsen et al.'s
    \cite{lauritsen2021framing} framing-instability claim on MIMIC-IV.

    \item \textbf{C5 (Equity-by-label):} To the best of our knowledge, the first
    study combining subgroup performance disparities with subgroup label
    sensitivity for any sepsis model, and the first to report both with
    cluster-bootstrap intervals under false-discovery control. The
    label-sensitivity disparity survives correction; the discrimination
    disparity does not, and reporting the second without the first would have
    inverted which finding is defensible.

    \item \textbf{C6 (Reproducible pipeline):} A pre-registered,
    TRIPOD+AI-conformant pipeline with fixed seeds and a single orchestration
    script, which reproduces every reported estimate exactly from the raw
    databases (Section~\ref{sec:code_availability}).

    \item \textbf{C7 (Definitional vs.\ empirical label variance separated):}
    A treatment-independent deterioration label, tested alongside the three
    Sepsis-3 variants, shows that the label effect is not an artefact of how one
    definition is read: it agrees with Sepsis-3 at $\kappa = 0.057$ and matches
    its onset hour in 6.4\% of shared stays. The same comparison shows that
    Sepsis-3 is harder to predict than either of the limbs it is built from,
    and that the pre-treatment window empty of Sepsis-3 onsets contains 1,362
    treatment-independent deterioration onsets identified at 0.840 AUROC. This
    relocates the difficulty of sepsis prediction from the physiology to the
    definition (Section~\ref{sec:label_anchor_results}).
\end{enumerate}


\section{Future Work}

\begin{enumerate}
    \item \textbf{Better sepsis labels.} Section~\ref{sec:label_anchor_results}
    takes the first step --- a deterioration label built from physiology alone,
    with no treatment timestamp --- and shows it identifies substantially
    different patients at substantially different hours. But that label is an
    operational construct built for this comparison, not a validated definition,
    and its components overlap the model's own inputs. The work remaining is a
    sepsis definition grounded in objective measurements such as lactate
    trajectories and organ-function biomarkers, adjudicated independently of the
    model's feature set, which would eliminate the treatment-constitutionality
    problem and allow genuine pre-onset prediction.

    \item \textbf{Non-Sepsis-3 consensus definitions.} SEP-1 and the CDC Adult
    Sepsis Event were not implemented. Neither would have tested
    treatment-independence --- both are defined on antibiotic days --- but each
    yields a different cohort, and Dutta et al.\ \cite{dutta2026performance}
    report that performance moves substantially between them. Adding them would
    extend the label-variance decomposition across definitions rather than
    within one.

    \item \textbf{Prospective validation.} Test the new label in real clinical
    settings where sepsis onset can be determined by a protocol-blinded assessor,
    independent of treatment decisions.

    \item \textbf{Utility-optimised thresholds.} Choose alert thresholds that
    maximise the Utility Score directly, rather than maximising AUROC. None of
    the evaluated models was optimised for clinical utility; a direct
    utility-maximising threshold search may recover positive utility at
    restricted operating points.

    \item \textbf{Multi-hospital evaluation.} The alert burden (over 90 false
    alerts per true alert under the best model) is a deployment problem, not a
    statistical one. Determining whether any model can be deployed at a
    practically acceptable burden requires prospective multi-centre testing with
    clinical-staff feedback.
\end{enumerate}


\section{Closing Remark}

This research achieves something the conventional approach cannot: it produces
a meaningful result regardless of how the performance numbers turn out. The
hypothesis framework ensures that confirmed, rejected, and superseded
hypotheses are all equally informative about the underlying problem.

The central finding, that the Sepsis-3 label on MIMIC-IV is constituted by
clinical action, is not a failure of the model. It is a finding about the label
itself, and it is the most important contribution this study makes.

The final chapter of that finding is the one that took a label outside Sepsis-3
to see. The pre-treatment window is not empty of predictable deterioration ---
it contains 1,362 onsets under a treatment-independent label, identified at
0.840 AUROC. It is empty of \emph{Sepsis-3}, because Sepsis-3 cannot place an
onset before the treatment it is anchored to. The signal clinicians would want a
model to find is present in the data. The definition the field has agreed to
train on is what cannot represent it.
